\documentclass{article}
\usepackage{icomp2026_conference,times}

\input{math_commands.tex}

\usepackage{hyperref}
\usepackage{url}
\usepackage{booktabs}
\usepackage{graphicx}
\usepackage{amsmath}
\usepackage{amssymb}

\title{DriftSSL: Self-Supervised Learning as Fisher Divergence Geometry}

\author{Anonymous Authors\\
Paper under double-blind review}

% Uncomment only for the camera-ready version after acceptance.
% \icompfinalcopy

\begin{document}
\maketitle

\begin{abstract}
We study contrastive self-supervised learning through a kernel-density view of representation geometry. The core object is a drift vector obtained from the difference between positive and negative kernel-weighted centroids, which gives a closed-form score-ratio field on the embedding sphere. This perspective links temperature, mean-shift dynamics, and Fisher divergence, and motivates diagnostics for training phases and representation collapse.
\end{abstract}

\section{Introduction}

Contrastive self-supervised learning is sensitive to temperature, augmentation strength, and representation geometry. We reinterpret this sensitivity through kernel density estimation on normalized embeddings.

\section{Method}

Let $z_i$ be a normalized embedding and let $\mu^+_\sigma(z_i)$ and $\mu^-_\sigma(z_i)$ denote kernel-weighted centroids of positive and negative neighborhoods. The drift vector is
$$
V_\sigma(z_i) = \mu^+_\sigma(z_i) - \mu^-_\sigma(z_i).
$$
Under a Gaussian kernel, $V_\sigma(z_i) / \sigma^2$ is the score of a local log-density ratio.

\section{Experiments}

Describe the datasets, SSL baselines, representation metrics, and downstream probes.

\section{Discussion}

Discuss what the drift-field view explains, where it is diagnostic rather than a direct loss, and how it relates to score matching and learning dynamics.

\bibliography{references}
\bibliographystyle{icomp2026_conference}

\appendix
\section{Additional Material}

Add derivations and extra experimental details here.

\end{document}
